% PATCH: R2-01 (6/12)
% Reviewer: 2
% Target section: sections/results.tex
% Requirement: Condense repetitive text — the six-frame compound-terrain sequence had
%   6 near-templated paragraphs ("In Frame N... the platform... mode... because..."),
%   each already echoed by its own figure-panel caption. Replaced the 6 paragraphs with
%   a compact table (Frame | Terrain segment | Mode | Trigger) plus one composite
%   paragraph tying the sequence together and preserving the key "validated across 15
%   trials" claim. Figure block (6-panel figure + its captions) left untouched.
% Status: applied

\begin{table}[htbp]
\centering
\caption{\revblue{Gait-mode transitions across the six-frame compound-terrain sequence (Figure~\ref{fig:real_variable_topography_sequence}), replicated across 15 continuous trials.}}
\label{tab:terrain_sequence}
\begin{tabularx}{\textwidth}{c l l X}
\toprule
\textbf{Frame} & \textbf{Terrain segment} & \textbf{Mode} & \textbf{Trigger / mechanism} \\
\midrule
1 & Double-ramp ascent & Quadrupedal walking & Static stability \& traction priority; elevated support polygon keeps CoM within bound during climb \\
2 & Transition to flat elevated segment & $\rightarrow$ Differential rolling & MLP-detected drop in vertical-IMU acceleration std.\ dev.\ below threshold suppresses quadrupedal salience, disinhibits rolling channel \\
3 & Horizontal elevated plane & Differential rolling (sustained) & Energy-conservative wheeled traversal; rolling channel remains stable \\
4 & Controlled ramp descent & Differential rolling (sustained) & Pitch disturbance stays below quadrupedal re-engagement threshold ($\approx 5^{\circ}$) \\
5 & Perlin-noise rough terrain & Rolling $\rightarrow$ stagnation & Wheels lose traction on high-variance surface; MLP stagnation channel triggers sustained signal \\
6 & Rough terrain (recovery) & $\rightarrow$ Quadrupedal walking & Stagnation triggers reactive disinhibition of quadrupedal channel; limbs redeploy, restoring support polygon \\
\bottomrule
\end{tabularx}
\end{table}

\revblue{Across the compound course (double-ascent ramp $\rightarrow$ elevated plane $\rightarrow$ descent ramp $\rightarrow$ Perlin-noise rough terrain, 15 trials), the arbitration network autonomously transitions through the gait sequence summarised in Table~\ref{tab:terrain_sequence} (Figure~\ref{fig:real_variable_topography_sequence}). The critical event occurs in Frames 5--6: rolling-mode wheel slip on the stochastic rough segment triggers a sustained stagnation signal, resolved by reactively redeploying the quadrupedal gait, which restores the support polygon and completes traversal. This closed-loop reactive switching was observed across all 15 trials, confirming reliable, autonomous failure detection and gait recovery under stochastic, non-structured terrain conditions.}
